\documentclass[11pt,a4paper]{article}
\usepackage[T1]{fontenc}
\usepackage{isabelle,isabellesym}

\usepackage{a4wide}
\usepackage{amssymb}

% this should be the last package used
\usepackage{pdfsetup}

% urls in roman style, theory text in math-similar italics
\urlstyle{rm}
\isabellestyle{it}

\begin{document}

\title{Soundness of the GKR Protocol Assembly}
\author{Jinwook Kim}

\maketitle

\begin{abstract}
The GKR protocol (Goldwasser, Kalai and Rothblum~\cite{gkr2015}) proves
the correct evaluation of a layered arithmetic circuit by reducing each
layer to a sumcheck instance and chaining the reductions down to the
circuit inputs.  This entry mechanizes the soundness of such an
assembly over the existing AFP formalization of the sumcheck protocol
by Garvia, Sprenger and Bootle~\cite{Sumcheck_Protocol-AFP}.

The development covers: a layered circuit intermediate representation
with a weighted gate family of algebraic degree at most three and its
evaluation semantics; multilinear extensions of value tables over the
boolean hypercube and the associated Lagrange weights; the per-layer
sumcheck identity connecting a layer's wiring predicates to the values
of the layer below; the consumption of the AFP sumcheck soundness and
completeness theorems for one layer reduction, via an explicit
multivariate-polynomial representative of the layer integrand; a
union-bound soundness theorem for the full layer-reduction chain,
including the two-point claim carry between layers; the equality of a
derived closed-form wiring oracle with the defining table oracle under
repartition completeness; and the correctness of proof batching as a
map over an abstract single-proof pipeline.

The protocol model is interactive and public-coin, with challenges
drawn uniformly from a finite field, matching the AFP sumcheck
formalization; the Fiat--Shamir transformation is outside the model.
The chain soundness bound is
$(s_{\mathit{out}} + \sum_i (2 s_i d + 1))/|\mathbb{F}|$ for layer
widths $2^{s_i}$ and a per-round degree bound $d$.  Related
mechanization efforts in other proof assistants include
ArkLib~\cite{arklib} in Lean; an expository account of the GKR
protocol and sumcheck-based proof systems is given by
Thaler~\cite{thaler2022}.
\end{abstract}

\tableofcontents

\newpage

% sane default for proof documents
\parindent 0pt\parskip 0.5ex

% generated text of all theories
\input{session}

% optional bibliography
\bibliographystyle{abbrv}
\bibliography{root}

\end{document}
